\section{Техническое задание}
\subsection{Основание для разработки}

Основанием для разработки является задание на выпускную квалификационную работу бакалавра "<Программное обеспечение для клавиатуры Charybdis 4x6 Mk2">

\subsection{Цель разработки}

Цель разработки -- создание единого удобного интерфейса для настройки сплит клавиатур без необходимости программирования -- позволить пользователям настраивать клавиатуру через графический интерфейс вместо правки кода прошивки.

\subsection{Назначение разработки}

Настройка раскладки клавиатуры:

\begin{itemize}
\item визуальное редактирование расположения клавиш;
\item создание и переключение между слоями;
\item поддержка модульной конструкции.
\end{itemize}

Программирование макросов:

\begin{itemize}
\item визуальный редактор макросов с возможностью записи нажатий;
\item поддержка сложных последовательностей;
\item назначение макросов на клавиши или комбинации.
\end{itemize}

Управление подсветкой:

\begin{itemize}
\item выбор из предустановленных эффектов;
\item настройка цвета отдельных клавиш или зон;
\item синхронизация подсветки между половинами клавиатуры;
\item создание пользовательских эффектов.
\end{itemize}

Работа с профилями:

\begin{itemize}
\item сохранение нескольких профилей настроек;
\item быстрое переключение между профилями;
\item экспорт/импорт профилей для обмена с другими пользователями.
\end{itemize}

Обновление прошивки:

\begin{itemize}
\item автоматическое обнаружение новых версий прошивки;
\item безопасное обновление без риска «окирпичивания» клавиатуры;
\item откат к предыдущей версии при необходимости.
\end{itemize}

Диагностика и отладка:

\begin{itemize}
\item мониторинг состояния клавиатуры;
\item тестирование клавиш на отклик;
\item логирование ошибок и проблем подключения.
\end{itemize}

Интеграция с ОС:

\begin{itemize}
\item поддержка горячих клавиш ОС;
\item интеграция с системными настройками;
\item уведомления о состоянии клавиатуры.
\end{itemize}

\subsection{Требования к функциональности}

В программном обеспечении должно быть представлено одно действующее лицо -- это пользователь.

На основании требований к программной системе в программе должны быть реализованы следующие функции:

\begin{enumerate}
\item \textbf{Работа с файлами.} Пользователь может сохранить текущую компоновку, включая все слои и макросы, в файл для последующего восстановления или переноса. Также поддерживается загрузка ранее сохраненной конфигурации.
\item \textbf{Языковая раскладка.} Пользователь может выбрать один из предложенных макетов.
\item \textbf{Безопасность.} Реализовать функцию блокировки клавиатуры, которая защищает от случайных или ненужных изменений.
\item \textbf{Тема.} Пользователи могут выбирать внешний вид программы, переключаясь между светлой и темной темами или настраивая цветовую схему в соответствии со своими предпочтениями.
\item \textbf{Раскладка клавиатуры.} Предоставьте инструменты для просмотра и изменения распределения каждой клавиши. Все изменения будут отображаться на виртуальной схеме устройства.
\item \textbf{Макрос.} Разрешает создавать, редактировать и назначать макросы на любые клавиши. Макросы могут содержать последовательности нажатий клавиш, задержки и другие команды.
\item \textbf{Подсветка.} Если клавиатура поддерживает подсветку, пользователь может контролировать цвет, яркость и визуальные эффекты клавиш.
\item \textbf{Степ (коснуться/долгое нажатие).} Разрешает назначать разные действия для короткого и долгого нажатия одной и той же кнопки, тем самым расширяя функциональность устройства.
\item \textbf{Комбинация.} Пользователь может создать свои собственные действия, которые активируются при одновременном нажатии нескольких клавиш. Например, комбинация клавиш A+S может имитировать нажатие ESC.
\item \textbf{Слои.} Система поддерживает работу с слоями. Пользователь может переключаться между ними и настраивать функции всей клавиатуры для каждого слоя отдельно.
\item \textbf{Тестирование матрицы клавиатуры.} Встроенный инструмент используется для проверки рабочего состояния всех клавиш. При нажатии на кнопку она будет подсвечиваться на схеме, что позволит быстро обнаружить неисправность.
\item \textbf{Переопределение ключей.} Разрешает настраивать модификаторы (Shift, Ctrl, Alt) и их комбинации с другими клавишами, изменяя стандартное поведение системы.
\item \textbf{Замена повторяющейся клавиши.} Предоставляет детальные настройки функции повторения последней нажатой клавиши или указания альтернативного действия при длительном нажатии, что очень удобно для навигации и редактирования текста.
\item \textbf{Основные настройки прошивки QMK.} Откройте доступ к параметрам расширения клавиатуры: настройте тайм-аут, измените поведение модификаторов и другие низкоуровневые параметры для оптимизации устройства.

\end{enumerate}

\subsection{Требования к интерфейсу}

Интерфейс представляет собой современное десктопное приложение с единым минималистичным стилем, ориентированное на интуитивно понятную навигацию для двух основных сценариев: пользовательского (просмотр и тестирование настроек) и административного (управление раскладкой, макросами, слоями и параметрами).

Основные элементы интерфейса:

\begin{enumerate}
\item Окно «Файл» позволяет пользователю сохранять существующие макеты, слои и макросы в файл или экспортировать этот файл для обработки.
\item Окно «Языковая раскладка» позволяет пользователю просматривать и изменять текущую раскладку клавиатуры.
\item Окно «Безопасность» позволяет пользователю заблокировать клавиатуру, чтобы предотвратить дальнейшие изменения.
\item Окно «Темы» позволяет пользователю изменять внешний вид программы.
\item Окно «Раскладка клавиатуры» позволяет пользователю просматривать или изменять распределение клавиш на текущей клавиатуре.
\item Окно "Макрос" позволяет пользователю создать макрос и записать его на определённую клавишу.
\item Окно "Подсветка" позволяет пользователю изменять подсветку клавиатуры, если такая функция поддерживается материнской платой.
\item Окно "Степ" позволяет пользователю настраивать различные действия в зависимости от способа использования кнопки. Например, одна и та же кнопка может быть настроена на активацию макроса при нажатии на кнопку или на активацию других функций при удерживании кнопки.
\item Окно "Комбинация" позволяет пользователям добавлять пользовательские действия при нажатии определенных комбинаций клавиш. Например, нажатие клавиш A и S в пределах комбинации приведет к нажатию клавиши ESC или выполнению более сложных задач.
\item Окно "Слой" позволяет пользователю изменять функции всей клавиатуры в зависимости от текущего "слоя" клавиатуры.
\item Окно «Переопределение ключей» позволяет пользователям проверять и настраивать поведение комбинаций модификаторов, переопределяя их реакцию при одновременном нажатии с любой другой клавишей.
\item Окно «Замена повторяющейся клавиши» позволяет пользователям детально настраивать повтор последней нажатой клавиши и повторно выполнять заменяющее действие. Этот инструмент идеально подходит для оптимизации взаимодействия с клавиатурой, создавая удобные и быстрые способы перемещения, редактирования текста и навигации.
\item Окно «Основные настройки прошивки QMK» позволяет пользователям точно настраивать рабочие параметры прошивки QMK, включая установку временных интервалов, настройку поведения модификаторов и т.д. Этот инструмент идеально подходит для оптимизации работы клавиатуры, повышения комфорта и продуктивности пользователя.
\item Окно "Тестирование матрицы клавиатуры" позволяет пользователям проверять работоспособность всех клавиш клавиатуры. Окно "Тестирование матрицы клавиатуры" позволяет пользователям проверять работоспособность всех клавиш клавиатуры.
\end{enumerate}

Этот метод удобен как для новичков, так и для опытных пользователей, и позволяет централизованно управлять всеми функциями Vial/QMK.

\subsection{Моделирование вариантов использования}

В системе выделяется единственное лицо – пользователь.

Далее приведены агрегированные перечни прецедентов для каждой роли. Они детализируются в последующих подпунктах.

Перечень прецедентов пользователя:

\begin{enumerate}
\item Прецедент "<Редактирование файла">.
\item Прецедент "<Смена языковой раскладки">.
\item Прецедент "<Настройки безопасности">.
\item Прецедент "<Смена темы">.
\item Прецедент "<Изменение раскладки клавиатуры">.
\item Прецедент "<Редактирование макроса">.
\item Прецедент "<Редактирование подсветки">.
\item Прецедент "<Изменение степа">.
\item Прецедент "<Редактирование комбинаций">.
\item Прецедент "<Изменение слоёв">.
\item Прецедент "<Переопределение ключей">.
\item Прецедент "<Изменение альтернативных клавиш повтора">.
\item Прецедент "<Смена базовых настроек прошивки QMK">. 
\item Прецедент "<Тестирование матрицы клавиатуры">.
\end{enumerate}

На рисунке \ref{usediagram:image} показана диаграмма вариантов использования программной системы.

\begin{figure}[H]
	\centering
	\includegraphics[width=1\linewidth]{usediagram}
	\caption{Диаграмма вариантов использования программной системы}
	\label{usediagram:image}
\end{figure}

\subsubsection{Сценарий для прецедента "<Редактирование файла">}

\textbf{Основной исполнитель:} пользователь

\textbf{Заинтересованные лица и их требования:} пользователь хочет импортировать или экспортировать настройки клавиатуры через файлы формата VIL.

\textbf{Предусловие:} необходимо, чтобы пользователь открыл программу и перешёл в раздел "<File">.

\textbf{Постусловие:} результат импорта или экспорта отображается на экране, либо сохраняется в файл.

\textbf{Основной успешный сценарий:}

\begin{enumerate}
\item Пользователь открывает программу и переходит в раздел "<File">.
\item Пользователь выбирает пункт "<Load saved Layout"> или "<Save current Layout">.
\item[3.1.] При выборе пункта "<Load saved Layout">:

\begin{itemize}
	\item пользователь загружает ранее созданный VIL-файл с настройками;
	\item программа импортирует настройки и отображает их на панели конфигурации.
\end{itemize}

\item[3.2.] При выборе пункта "<Save current Layout">:

\begin{itemize}
	\item пользователь сохраняет текущие настройки в новый VIL-файл;
	\item файл сохраняется на компьютер пользователя.
\end{itemize}

\end{enumerate}

\subsubsection{Сценарий для прецедента "<Смена языковой раскладки">}

\textbf{Основной исполнитель:} пользователь

\textbf{Заинтересованные лица и их требования:} пользователь хочет изменить языковую раскладку на другую из предложенного ему списка.

\textbf{Предусловие:} необходимо, чтобы пользователь открыл программу и перешел в раздел "<Keyboard Layout">.

\textbf{Постусловие:} результат применения языковой раскладки отображается на экране программы или сохраняется локально.

\textbf{Основной успешный сценарий:}

\begin{enumerate}
\item Пользователь запускает программу и переходит в раздел "<Keyboard Layout">.
\item Пользователь выбирает один из вариантов языковой раскладки.
\item После завершения процесса пользователю на экран выводится та языковая раскладка, которую он выбрал.
\end{enumerate}

\subsubsection{Сценарий для прецедента "<Настройки безопасности">}

\textbf{Основной исполнитель:} пользователь

\textbf{Заинтересованные лица и их требования:} пользователь хочет разблокировать клавиатуру для изменения в ней раскладки, настроек клавиш.

\textbf{Предусловие:} необходимо, чтобы пользователь открыл программу и перешел в раздел "<Security">.

\textbf{Постусловие:} клавиатура будет заблокирована, либо разблокирована.

\textbf{Основной успешный сценарий:}

\begin{enumerate}
\item Пользователь запускает программу и переходит в раздел "<Security">.
\item Пользователь выбирает одну из операций:

\begin{itemize}
	\item разблокировать ("<Unlock">);
	\item заблокировать ("<Lock">);
\end{itemize}

\item[3.1.] Разблокировать ("<Unlock">): разблокирует клавиатуру для изменений при зажатии определённых клавиш на определённое время.
\item[3.2.] Заблокировать ("<Lock">): блокирует клавиатуру для изменений.
\end{enumerate}

\subsubsection{Сценарий для прецедента "<Смена темы">}

\textbf{Основной исполнитель:} пользователь

\textbf{Заинтересованные лица и их требования:} пользователь хочет изменить внешний вид интерфейса программы через список предоставленных ему внешних видов.

\textbf{Предусловие:} необходимо, чтобы пользователь открыл программу и перешел в раздел "<Theme">.

\textbf{Постусловие:} выбранная тема отображается на экране.

\textbf{Основной успешный сценарий:}

\begin{enumerate}
\item Пользователь запускает программу и переходит в раздел "<Theme">.
\item Пользователь выбирает один из вариантов тем, который удовлетворяет.
\item По завершении операции, пользователю выводится сообщение о том, что выбранная им тема была применена в программе.
\end{enumerate}

\subsubsection{Сценарий для прецедента "<Изменение раскладки клавиатуры">}

\textbf{Основной исполнитель:} пользователь

\textbf{Заинтересованные лица и их требования:} пользователь хочет импортировать или экспортировать карту клавиш (keymap) клавиатуры через файлы формата VIL.

\textbf{Предусловие:} необходимо, чтобы пользователь открыл программу и перешел в раздел "<Keymap">.

\textbf{Постусловие:} результат импорта или экспорта отображается на панели конфигурации или сохраняется локально.

\textbf{Основной успешный сценарий:}

\begin{enumerate}
\item Пользователь запускает программу и переходит в раздел "<Keymap">.
\item Пользователь выбирает один из пунктов:

\begin{itemize}
	\item импорт карты клавиш ("<Import Keymap">);
	\item экспорт карты клавиш ("<Export Keymap">).
\end{itemize}

\item[3.1.] Импорт карты клавиш:

\begin{itemize}
	\item пользователь загружает ранее созданный файл формата VIL (*.vil) c настройками keymap;
	\item программа импортирует указанные настройки и отображает обновленную карту клавиш на панели конфигурации.
\end{itemize}

\item[3.2.] Экспорт карты клавиш:

\begin{itemize}
	\item пользователь инициирует процесс сохранения текущей карты клавиш;
	\item программа создает новый файл формата VIL (*.vil) с настройками keymap, и сохраняет его на устройство пользователя.
\end{itemize}

\item[4.] По окончании процедуры пользователю показывается готовая раскладка клавиатуры.
\end{enumerate}

\subsubsection{Сценарий для прецедента "<Редактирование макроса">}

\textbf{Основной исполнитель:} пользователь

\textbf{Заинтересованные лица и их требования:} пользователь хочет импортировать или экспортировать макросы клавиатуры через файлы формата VIL.

\textbf{Предусловие:} необходимо, чтобы пользователь открыл программу и перешел в раздел "<Macros">.

\textbf{Постусловие:} макросы успешно импортируются или экспортируются, и результат отображается на панели конфигурации или сохраняется локально.

\textbf{Основной успешный сценарий:}

\begin{enumerate}
\item Пользователь запускает программу и переходит в раздел "<Macros">.
\item Пользователь выбирает один из пунктов:

\begin{itemize}
	\item открыть текстовый редактор (Open Text Editor);
	\item добавить действие (Add action);
	\item нажать клавишу «Enter» (Tap Enter);
	\item записать макрос (Record macro).
\end{itemize}

\item[3.1.] Открыть текстовый редактор:

\begin{enumerate}
	\item Открывается консоль, в которой можно прописать макрос. Макрос – это список, каждый элемент которого это действие. Каждое действие – это список, где первый элемент – это название действия, последующие элементы – это аргументы действия. Пример:
	
	\begin{figure}[H]
		\begin{lstlisting}[language=Python]
		[["down", "KC_W", "KC_T", "KC_U"], ["up", "KC_W", "KC_T", "KC_U"],
		["delay", 42000], ["down", "KC_Z"], ["up", "KC_Z"], ["delay", 15],
		["tap", "KC_U"], ["text", "Primer"]]
		\end{lstlisting}
	\end{figure}
	\item После чего можно как применить те строки кода, которые вы написали, и они автоматически построят алгоритм макроса.
\end{enumerate}

\item[3.2.] Добавить действие:

При нажатии на данную кнопку в поле для создания макроса появляется строка, в которой можно поменять действие на то, что предлагается в меню:

\begin{itemize}
	\item нажать клавишу (Down);
	\item отжать клавишу (Up);
	\item один раз нажать клавишу (Tap);
	\item ввести заданный текст (Text).
\end{itemize}

\item[3.3.] Нажать клавишу «Enter»:

При нажатии клавиши добавляется в строку действие нажатия клавиши «Enter» в панель редактирования макроса.

\item[3.4.] Записать макрос:

При нажатии открывается панель с выбором двух действий:

\begin{itemize}
	\item добавить к текущему макросу (Append to current);
	\item изменить текущий макрос (Replace everything). 
\end{itemize}

\item[4.] По завершении процедуры написания макроса его можно сохранить и добавить на конкретную клавишу на клавиатуре.
\end{enumerate}

\subsubsection{Сценарий для прецедента "<Редактирование подсветки">}

\textbf{Основной исполнитель:} пользователь

\textbf{Заинтересованные лица и их требования:} пользователь хочет изменить настройки подсветки клавиатуры.

\textbf{Предусловие:} необходимо, чтобы пользователь открыл программу и перешел в раздел "<Lighting">.

\textbf{Постусловие:} результат изменения подсветки отображается на панели конфигурации и на самом устройстве.

\textbf{Основной успешный сценарий:}

\begin{enumerate}
\item Пользователь запускает программу и переходит в раздел "<Lighting">.
\item Пользователь настраивает следующие показатели:

\begin{enumerate}
	\item Эффект подсветки (RGB Effect): открывается список возможных алгоритмов подсветки.
	\item Цвет подсветки (RGB Color): открывается типичная панель для выбора конкретного цвета подсветки.
	\item Яркость подсветки (RGB Brightness): тут в зависимости от положения ползунка меняется яркость, с которой работает алгоритм подсветки.
	\item Скорость подсветки (RGB Speed): тут в зависимости от положения ползунка меняется скорость, с которой работает алгоритм подсветки.
\end{enumerate}

\item[4.] При нажатии на кнопку сохранить (save) настройки автоматически применяются.
\end{enumerate}

\subsubsection{Сценарий для прецедента "<Изменение степа">}

\textbf{Основной исполнитель:} пользователь

\textbf{Заинтересованные лица и их требования:} пользователь хочет импортировать или экспортировать последовательности действий тап-данса клавиатуры через файлы формата VIL.

\textbf{Предусловие:} необходимо, чтобы пользователь открыл программу и перешел в раздел "<Tap Dance">.

\textbf{Постусловие:} последовательности степа успешно импортируются или экспортируются, и результат отображается на панели конфигурации или сохраняется локально.

\textbf{Основной успешный сценарий:}

\begin{enumerate}
\item Пользователь запускает программу и переходит в раздел "<Tap Dance">.
\item[2.1. ] Пользователь загружает готовый файл клавиатуры, в котором сохранены предварительно написанные «Степы».
\item[2.2. ] Пользователь пишет сам: тут уже пользователь выбирает разные настройки, которые ему нужны и клавиши, которые должны прожиматься: при нажатии (On tap), при удержании (On hold), при двойном нажатии (Double tap) при нажатии и удержании (On tap + hold). Также в этих параметрах можно задать время в миллисекундах, которое должно пройти, чтоб клавиатура среагировала на «Степ». 
\item[3. ] Импорт последовательностей тап-данса:

\begin{enumerate}
	\item[3.1. ] Выбрать нужную клавишу в нужном слое клавиатуры, на которой будет установлен «Степ».
	\item[3.2. ] Выбрать из списка тот «Степ», который нужен.
\end{enumerate}

\end{enumerate}

\subsubsection{Сценарий для прецедента "<Редактирование комбинаций">}

\textbf{Основной исполнитель:} пользователь

\textbf{Заинтересованные лица и их требования:} пользователь хочет отредактировать или создать новые комбинации клавиш.

\textbf{Предусловие:} необходимо, чтобы пользователь открыл программу и перешел в раздел "<Combos">.

\textbf{Постусловие:} комбинации клавиш успешно изменяются или добавляются, и результат отображается на панели конфигурации или сохраняется локально.

\textbf{Основной успешный сценарий:}

\begin{enumerate}
\item Пользователь запускает программу и переходит в раздел "<Combos">.
\item[2.1. ] Пользователь импортирует комбинации клавиш из сохранённого файла:
\item[2.1.1. ] Пользователь загружает ранее созданный файл формата VIL (*.vil), содержащий комбо-настройки.
\item[2.1.2. ] Программа импортирует комбинации и отображает их на панели конфигурации.
\item[2.2. ] Пользователь cоздаёт своё комбо:
\item[2.2.1. ] Пользователь выбирает до 4-х клавиш, которые могут быть прожаты одновременно.
\item[2.2.2. ] Пользователь выбирает клавишу, которая будет на выходе.
\item[3. ] По завершении процедуры пользователь может пользоваться только что выгруженным комбо, либо только что сделанным с нуля.
\end{enumerate}

\subsubsection{Сценарий для прецедента "<Изменение слоёв">}

\textbf{Основной исполнитель:} пользователь

\textbf{Заинтересованные лица и их требования:} пользователь хочет импортировать, или создать, или отредактировать оригинальные слои клавиатуры через файл формата VIL.

\textbf{Предусловие:} необходимо, чтобы пользователь открыл программу и перешел в раздел "<Layers">.

\textbf{Постусловие:} слои клавиатурных раскладок успешно импортируются, или будут созданы, или будут отредактированы, и результат отображается на панели конфигурации и сохраняется локально.

\textbf{Основной успешный сценарий:}

\begin{enumerate}
\item Пользователь запускает программу и переходит в раздел "<Layers">.
\item[2.1. ] Пользователь загружает слой: выгрузка слоя происходит вместе с импортом исходного файла с разрешением VIL.
\item[2.2. ] Пользователь редактирует слой:

\begin{enumerate}
	\item Пользователю следует нажать на номер того слоя, который он хочет отредактировать (нумерация идёт с нулевого слоя).
	\item Пользователю следует выбрать конкретную клавишу в слое.
	\item Пользователю следует выбрать из предложенного перечня функций те, которые ему нужны в количестве 1 штука на конкретно выбранную клавишу.
\end{enumerate}

\item[2.3. ] Пользователь создаёт слой: создание происходит аналогично редактированию, только редактировать нужно все кнопки внутри слоя вручную.
\item[3. ] По завершении процедуры пользователь либо выгружает готовые слои из файла, либо отредактирует нужный слой, либо создаст свой собственный, удобный непосредственно для негой слой.
\end{enumerate}

\subsubsection{Сценарий для прецедента "<Переопределение ключей">}

\textbf{Основной исполнитель:} пользователь

\textbf{Заинтересованные лица и их требования:} пользователь хочет импортировать или создать переопределения клавиш (key overrides) клавиатуры через файлы формата VIL.

\textbf{Предусловие:} необходимо, чтобы пользователь открыл программу и перешел в раздел "<Key Overrides">.

\textbf{Постусловие:} переопределения клавиш успешно импортируются или создаются, и результат отображается на панели конфигурации или сохраняется локально.

\textbf{Основной успешный сценарий:}

\begin{enumerate}
\item Пользователь запускает программу и переходит в раздел "<Key Overrides">.
\item[2.1. ] Пользователь импортирует переопределение клавиш: импортируются все переопределённые клавиши при импортировании файла раскладки.
\item[2.2. ] Пользователь создаёт переопределение клавиш: выбирает нужные для него параметры и клавишу, из-за которой сработают выбранные параметры.
\item[3.] По завершении процедуры пользователь, либо загрузит себе переопределённые клавиши, либо создаст их вручную.
\end{enumerate}

\subsubsection{Сценарий для прецедента "<Изменение альтернативных клавиш повтора">}

\textbf{Основной исполнитель:} пользователь

\textbf{Заинтересованные лица и их требования:} пользователь хочет включить, или отключить, или задать своё сочетание альтернативных клавиш. 

\textbf{Предусловие:} необходимо, чтобы пользователь открыл программу и перешел в раздел "<Alt Repeat Key">.

\textbf{Постусловие:} настройки поведения повторения клавиш при удерживании модификатора Alt отображаются на панели конфигурации или сохраняются локально.

\textbf{Основной успешный сценарий:}

\begin{enumerate}
\item Пользователь запускает программу и переходит в раздел "<Alt Repeat Key">.
\item[2.1. ] Пользователь включает сочетание альтернативных клавиш: ставит флажок включить.
\item[2.2. ] Пользователь отключает сочетание альтернативных клавиш: убирает флажок включить.
\item[2.3. ] Пользователь задаёт сочетание альтернативных клавиш:

\begin{itemize}
	\item добавляет предыдущую клавишу;
	\item добавляет альтернативную;
	\item может добавить какие-то другие клавиши, представленные из списка, чтоб эта альтернативная клавиша сработала;
	\item может выбрать одну из настроек.
\end{itemize}

\item[3. ] После успешного завершения операции у пользователя автоматически сохраняются все изменения в соответствующий файл с разрешением VIL.
\end{enumerate}

\subsubsection{Сценарий для прецедента "<Смена базовых настроек прошивки QMK">. }

\textbf{Основной исполнитель:} пользователь

\textbf{Заинтересованные лица и их требования:} пользователь хочет изменить  настройки прошивки QMK клавиатуры.

\textbf{Предусловие:} необходимо, чтобы пользователь открыл программу и перешел в раздел "<QMK Settings">.

\textbf{Постусловие:} настройки прошивки QMK успешно изменены, и результат отображается на панели конфигурации и сохраняется локально.

\textbf{Основной успешный сценарий:}

\begin{enumerate}
\item Пользователь запускает программу и переходит в раздел "<QMK Settings">.
\item Пользователь включает или выключает те пункты, которые ему нужны.
\end{enumerate}

\subsubsection{Сценарий для прецедента "<Тестирование матрицы клавиатуры">}

\textbf{Основной исполнитель:} пользователь

\textbf{Заинтересованные лица и их требования:} пользователь хочет провести тестирование и диагностику оборудования с помощью инструмента «Matrox Tester».

\textbf{Предусловие:} необходимо, чтобы пользователь открыл программу и перешел в раздел "<Matrox Tester">.

\textbf{Постусловие:} на экране отображаются все нажатые и не нажатые клавиши.

\textbf{Основной успешный сценарий:}

\begin{enumerate}
\item Пользователь запускает программу и переходит в раздел "<Matrox Tester">.
\item Пользователь нажимает на клавиши.
\item На экране нажатые клавиши подсвечиваются серым.
\end{enumerate}

\subsection{Требования к надёжности}

Система обеспечивает стабильную работу и защиту пользовательских настроек как в штатном режиме, так и при возникновении ошибок.

\textbf{Целостность данных.} Все операции по созданию и изменению раскладок, макросов, слоёв и профилей выполняются атомарно. Это гарантирует, что в памяти контроллера не сохранятся частично записанные или повреждённые конфигурации. При сбое во время записи система откатывает изменения, возвращая клавиатуру к последнему рабочему состоянию.

\textbf{Устойчивость работы.} Прошивка и сопутствующее ПО корректно обрабатывают внутренние ошибки. В случае некорректных действий пользователя или аппаратных сбоев система не "<зависает"> и не требует сброса питания, а выводит понятное уведомление и предлагает варианты решения.

\textbf{Валидация данных.} На всех этапах -- от ввода в интерфейсе до записи в память клавиатуры -- выполняется строгая проверка данных. Система не позволяет сохранить пустые макросы, некорректные назначения клавиш или нарушенные комбинации. Пользователь получает мгновенную обратную связь при попытке ввести недопустимые значения.

\textbf{Логирование событий.} Все важные действия и ошибки фиксируются в журнале. В случае сбоя пользователь или разработчик может просмотреть лог, чтобы понять причину проблемы. Записывается контекст: какое действие выполнялось, какие параметры были заданы и на каком этапе возникла ошибка.

\textbf{Резервное копирование.} Для защиты от потери всех настроек реализован механизм резервного копирования. Пользователь может в любой момент сохранить полный образ конфигурации (прошивки) в файл на компьютере. Рекомендуется регулярно создавать такие бэкапы, особенно перед глобальными изменениями или обновлением прошивки. Восстановление из резервной копии происходит быстро и безопасно.

\subsection{Требования к информационной безопасности}

Система обеспечивает многоуровневую защиту от несанкционированного доступа, случайных ошибок и вредоносных действий, гарантируя безопасность пользовательских настроек и стабильность работы устройства.

\textbf{Разграничение доступа к административным функциям.} Доступ к разделу управления (эквивалент "<админки"> в Vial Configurator) строго ограничен. Вносить изменения в раскладку, макросы, слои и системные параметры может только авторизованный пользователь. Это предотвращает случайное или преднамеренное изменение конфигурации посторонними лицами, имеющими физический или сетевой доступ к компьютеру.

\textbf{Защита от некорректных и вредоносных данных.} Все данные, передаваемые от редактора к клавиатуре, проходят строгую проверку. Система не позволяет отправить на устройство некорректные команды, повреждённые файлы прошивки или данные, не соответствующие ожидаемому формату. Это защищает контроллер клавиатуры от "<окирпичивания"> (состояния, когда устройство перестаёт работать) и гарантирует, что загруженная конфигурация будет стабильной.

\textbf{Валидация пользовательского ввода.} При создании макросов, настройке комбинаций клавиш или изменении назначения кнопок система проверяет вводимые данные. Она контролирует длину последовательностей, корректность используемых клавиш и отсутствие запрещённых символов. Это предотвращает создание неработающих или вызывающих сбои команд.

\textbf{Защита критических операций.} Действия, которые могут привести к потере данных или изменению поведения клавиатуры (например, сброс к заводским настройкам, удаление всех макросов, загрузка новой прошивки), требуют дополнительного подтверждения от пользователя.

\subsection{Требования к программной и технической совместимости}

Программное обеспечение для настройки сплит-клавиатур спроектировано для максимально широкой совместимости и стабильной работы в самых распространённых пользовательских сценариях.

\textbf{Операционные системы.} Приложение корректно функционирует на всех современных десктопных платформах, включая Windows, macOS и Linux. Это обеспечивает свободу выбора для пользователя, независимо от его предпочтений в операционной системе.

\textbf{Программные зависимости.} Для работы редактора не требуется установка сложного серверного окружения. Основные требования минимальны:

\begin{enumerate}
\item Операционная система: Windows 10/11, macOS (последние версии) или современный дистрибутив Linux.
\item Среда выполнения: наличие установленного Python (версии 3.7 и выше) или Node.js, так как многие инструменты для работы с прошивками основаны на этих языках.
\item Драйверы: для корректного определения клавиатуры в системе должны быть установлены необходимые драйверы USB-устройств.
\end{enumerate}

\textbf{Аппаратная совместимость.} Ключевым требованием является поддержка клавиатурой протокола QMK. Приложение должно без проблем работать с большинством популярных сплит-клавиатур, представленных на рынке, при условии, что на них установлена совместимая прошивка. Соединение с устройством происходит по стандартному протоколу USB-HID.

\textbf{Совместимость данных.} Система обеспечивает полную переносимость пользовательских настроек. Все конфигурации -- раскладки, макросы, слои и профили -- хранятся в едином, кроссплатформенном формате (обычно это текстовые файлы в формате JSON или конфигурационные файлы прошивки). Это позволяет:

\begin{itemize}
\item легко переносить свои настройки между разными компьютерами;
\item делиться конфигурациями с другими пользователями;
\item хранить резервные копии в облачных сервисах.
\end{itemize}

\subsection{Требования к оформлению документации}

Разработка программной документации и программного изделия должна производиться согласно ГОСТ 19.102-77 и ГОСТ 34.601-90. Единая система программной документации.
